\chapter{Discussão}
\label{chap:discussao}

\section{Resposta à Questão Principal}

A questão principal deste trabalho pergunta como a eficácia das estratégias
de retreinamento varia em função do tipo e da severidade do drift e do
classificador monitorado. Os resultados indicam que as três variáveis
importam: as interações da estratégia com o tipo, com a severidade e com o
classificador foram todas significativas e com efeitos de magnitude
relevante (Tabela~\ref{tab:anova}). Essas interações, no entanto, alteraram
menos a ordenação entre as estratégias do que a distância entre elas. Na
maior parte das condições, a janela deslizante foi a estratégia mais eficaz,
e as variáveis do drift determinaram sobretudo quanto as estratégias
acionadas por detecção ficaram atrás dela.

O fator que melhor explica esse padrão é a política de acionamento, e não o
procedimento de ajuste do modelo. A janela deslizante e o reset completo
treinam o modelo exatamente da mesma forma e diferem apenas em quando o
fazem (Seção~\ref{sec:estrategias}). Sob drift abrupto, em que o HDDDM
detecta a mudança no próprio bloco em que ela ocorre, as duas estratégias
praticamente empataram. Sob drift gradual e incremental, o detector dispara
no meio da transição e, na maioria das vezes, uma única vez; o reset produz
então um modelo ajustado a uma distribuição intermediária, que permanece em
uso depois que a transição termina. A análise com detector ideal
(Tabela~\ref{tab:detector-ideal}) confirma essa leitura: quando acionado no
momento certo, o reset superou a janela nos padrões abrupto e recorrente,
mas continuou atrás dela nos padrões lentos, porque nenhuma estratégia que
atua uma única vez acompanha uma distribuição que continua a se deslocar.

Esse resultado qualifica a premissa, comum na literatura, de que a
adaptação informada por detecção é preferível à adaptação cega
\cite{gama2014}. Nas condições estudadas, a adaptação cega foi mais robusta
porque não depende de o detector perceber cada etapa da mudança. Essa
robustez, porém, não é gratuita nem incondicional: a janela retreinou cerca
de quatro vezes mais que o reset (Tabela~\ref{tab:custo}), e a análise de
sensibilidade mostrou que, com cadência de dezesseis blocos, ela perdeu para
o reset (Tabela~\ref{tab:sensibilidade}). A conclusão mais precisa é,
portanto, que uma janela com cadência suficientemente frequente em relação à
duração das transições supera estratégias acionadas por um detector que
dispara uma única vez por mudança, ao custo de mais retreinamentos.

\subsection{Tipo de drift}

Para o drift recorrente, a literatura sugere que reaproveitar modelos
históricos é mais eficiente que retreinar do zero \cite{sun2018,li2021}. O
mecanismo de recorrência funcionou como esperado: o repositório identificou
o modelo correto em todos os retornos detectados. Ainda assim, a reativação
não superou o reset em acurácia. A explicação está na construção do
experimento: todo modelo é treinado sobre um único bloco de 400 instâncias,
e um modelo antigo treinado sobre um bloco de $P$ não é melhor do que um
modelo novo treinado sobre outro bloco de $P$. A vantagem da reativação
aparece no custo, já que ela evita treinos (1,2 treinos por stream, contra
2,2 do reset), e deveria aparecer também em acurácia em cenários nos quais o
modelo histórico tenha sido treinado com muito mais dados do que os
disponíveis logo após a mudança, ou nos quais os rótulos cheguem com atraso,
situação em que \citeonline{rahmani2023} mostram que a disponibilidade de
rótulos condiciona o desempenho do retreinamento. Essa interpretação é uma
hipótese a ser testada, e não um resultado deste trabalho.

\subsection{Severidade}

A severidade foi o segundo fator de maior efeito e interagiu com a
estratégia, em concordância com a observação de \citeonline{wang2022elastic}
de que a magnitude do drift condiciona a readaptação do modelo. Seu efeito
mais visível foi ampliar o ganho de todas as estratégias sobre o controle e
a desvantagem do retreinamento incremental, cujo modelo preserva um
conhecimento tanto mais prejudicial quanto maior o deslocamento. O
reset completo só superou a janela em combinações de severidade alta ou média
com mudanças abruptas, nas quais descartar o modelo antigo de uma só vez é
precisamente o que a situação exige. Não se observou, por outro lado, uma
vantagem de estratégias conservadoras sob severidade baixa: o incremental
foi pior que o reset também nesse nível.

\subsection{Classificador}

A interação entre classificador e estratégia decorreu quase inteiramente do
retreinamento incremental, cuja eficácia dependeu de como cada classificador
incorpora dados novos. No XGBoost, as novas rodadas de boosting são
treinadas para corrigir os erros do conjunto existente sobre os dados
recentes, o que tornou o incremental equivalente ao reset. No Random Forest,
as árvores novas são minoria diante das árvores treinadas antes do drift, e
a média das predições permanece dominada pelo conceito antigo. No Naive
Bayes, as estatísticas acumuladas misturam as duas distribuições. Esse
resultado reforça, em um desenho controlado, a observação de
\citeonline{webb2016} e \citeonline{wang2022dynamic} de que um mesmo
mecanismo de adaptação produz efeitos distintos em classificadores
distintos: a escolha entre atualizar e substituir um modelo não pode ser
feita independentemente do classificador em uso.

\section{Resposta à Questão Derivada}

A questão derivada pergunta se sinais observáveis do detector permitem
esboçar uma orientação de escolha de estratégia. A resposta tem duas partes.
Os três sinais mostraram-se informativos: a severidade observada acompanhou
a severidade real, a velocidade da divergência separou mudanças abruptas de
lentas, e a distância ao repositório identificou todas as recorrências
detectadas. Esses resultados sustentam, em ambiente controlado, a proposta
de usar a própria estatística do detector como aproximação da severidade
\cite{komorniczak2022} e indicam que parte do entendimento do drift, o
componente que \citeonline{lu2018} apontam como o menos desenvolvido da
área, pode ser extraída de um detector baseado em distribuição sem rótulos.

A orientação construída a partir desses sinais, porém, não superou a
política fixa de sempre usar a janela deslizante. Isso não decorre de os
sinais serem pouco confiáveis, mas de a janela ter sido a melhor estratégia
em quase todas as condições: quando uma única estratégia domina, conhecer o
tipo e a severidade do drift acrescenta pouco à decisão, e até o oráculo,
que conhece ambos, superou a janela fixa por apenas 0,002. Os sinais tendem
a ser mais úteis quando o custo também entra na decisão. Nas condições em
que o reset empata com a janela, como o drift abrupto e o recorrente, os
sinais permitiriam optar por uma estratégia com cerca de um quarto dos
retreinamentos, sem perda relevante de acurácia; nas transições lentas,
indicariam que a janela compensa seu custo adicional. Uma orientação
sensível ao custo, e não apenas à acurácia, é o desdobramento mais direto
deste resultado.

\section{Implicações para Pipelines de MLOps}

Três implicações práticas decorrem dos resultados, ressalvados os limites
discutidos na seção seguinte. A primeira é que acionar o retreinamento
apenas pelo detector é arriscado quando a mudança é lenta: um pipeline que
dispara um único retreinamento por detecção tende a produzir um modelo
ajustado a uma distribuição de transição. Combinar o acionamento por
detecção com retreinamentos agendados, ou continuar retreinando após uma
detecção até que a divergência se estabilize, são alternativas que os
resultados com detector ideal e com a janela deslizante sugerem. A segunda é
que a escolha entre atualizar e substituir um modelo deve considerar o
classificador, já que a atualização incremental foi adequada para o XGBoost
e inadequada para o Random Forest. A terceira é que a estatística do
detector, já calculada para fins de monitoramento, carrega informação útil
sobre a severidade, a velocidade e a recorrência do drift, e pode ser
registrada junto com cada alerta a custo praticamente nulo, oferecendo aos
critérios de resposta que \citeonline{kreuzberger2023} apontam como ausentes
nos pipelines atuais uma base observável.

\section{Ameaças à Validade}
\label{sec:ameacas}

\textbf{Validade externa.} Os dados são sintéticos: quatro atributos
gaussianos independentes, classe binária equilibrada, drift aplicado sempre
a dois atributos e transições de duração fixa. A independência entre
atributos é também o que permite expressar a severidade observada como
$1 - \prod_k \widehat{BC}_k$. Os resultados não se transferem
automaticamente para dados reais, com atributos correlacionados, desbalanço
de classes e transições de duração variável; o estudo de base
\cite{rocha2026}, por exemplo, usa conjuntos de dados reais e sintéticos de
referência que não foram considerados aqui. Além disso, um único detector
foi avaliado, e outros detectores podem disparar de forma diferente nas
transições lentas.

\textbf{Validade interna.} Várias escolhas de projeto afetam os resultados:
o tamanho do bloco, o uso de um único bloco como janela de treino, a cadência
da janela deslizante e o parâmetro $\gamma$ do HDDDM. As duas últimas foram
avaliadas na análise de sensibilidade, que mostrou que a vantagem da janela
depende da cadência. A agenda da janela teve fase sorteada por stream para
evitar que coincidisse com o início do drift. Como o drift começa sempre no
início de um bloco, a latência de detecção nula observada nos padrões abrupto
e recorrente é otimista em relação a um cenário em que a mudança ocorre no
meio de um bloco.

\textbf{Validade de construto.} No gerador adotado, o deslocamento de $p(X)$
é sempre acompanhado de uma mudança em $p(y \mid X)$, de modo que cada drift
combina componentes reais e virtuais, e o experimento não permite separar o
efeito de cada componente. A eficácia foi medida por acurácia, adequada a
classes equilibradas, e o critério de recuperação (3 pontos percentuais da
acurácia pré-drift) é uma convenção; por isso, o tempo de recuperação é
apresentado junto com a fração de transições recuperadas.

\textbf{Validade de conclusão.} As quinze execuções de cada repetição
compartilham a mesma stream, o que foi tratado com uma ANOVA em parcelas
subdivididas e com testes pareados por stream. A avaliação da orientação
preliminar usou 60 streams de teste e classifica cada stream a partir de
todos os seus episódios, o que a torna uma avaliação retrospectiva do
mapeamento entre sinais e estratégias, e não de uma política executada em
tempo real.
